\clearpage
\section{Literature context and interpretation}
\label{sec:literature}
The archived frozen test identifies limits of the earlier FFT2D5 closures for a
specified source and observable. It does not show that all 2.5D methods fail.
Numerical convergence, physical-model accuracy and tracking accuracy remain
separate questions.

\subsection{The longitudinal closure matters}
The label 2.5D does not specify a unique solver. For example, ImpactX
distinguishes line-density-weighted transverse fields with a density-derivative
longitudinal kick from its Gaussian potential-based model and full 3D
models~\cite{impactx}. The archived OPALX operator solves separate transverse slices and adds
the open-space local closure
\begin{equation}
 E_s^{\rm local}(s)=-\frac{g_0}{4\pi\eps\gamma^2}
 \frac{\dd\lambda}{\dd s},\qquad g_0=6.36.
 \label{eq:literature-local}
\end{equation}
The value and formula are verified in the tested source code. For the
constant-section source, $\lambda$ is constant inside its sharp ends.
Therefore no finite choice of $g_0$ can recover a nonzero interior end field
where $\lambda'=0$. Refinement can resolve the density step more sharply
without recovering the finite-width Coulomb response in \cref{fig:frozen-end}.
This is our closure-level inference, independent of numerical end padding.

de Hoon et al.'s author manuscript derives a nonlocal Fourier--Bessel field and recovers the local
$g$ approximation by Taylor expansion; they explicitly identify end-region
breakdown~\cite{dehoon}. Their field is transversely averaged for a straight,
circularized beam in a conducting circular pipe. It supports the scope
diagnosis, not a quantitative prediction for our open elliptical arc; their
$R$ denotes pipe radius, not bending radius.

Oeftiger and Hegglin's equivalent longitudinal model uses a parabolic line
density; their Figure 1 gives $g_0=6.36$ for one free-space example. Their
$z_m>3r_p$ pipe-validity condition concerns a round conducting pipe,
not our free-space source~\cite{oeftiger}.
Oeftiger's thesis makes the ellipsoid's free-space $g_0$ shape-dependent~\cite{oeftigerthesis}.
Thus the matching code constant is a provenance clue, not a universal value
or proof of the original developer's intent.

Baartman distinguishes on-axis from transversely averaged forces in a round
grounded pipe; even their additive constants differ~\cite{baartman}.
Pipe, open-space, averaging and transverse-profile conventions must therefore
be matched before comparing numerical $g$ values.

Venturini and Wang--Li retain wavelength-dependent longitudinal response
rather than one constant~\cite{venturini,wangli}. Wang--Li include a 20 keV
$\mathrm H_2^+$ case with $\beta\simeq0.0046$ and 5 mm radius, but for a
straight round beam. These kernels motivate a nonlocal longitudinal model; they are not
already references for our finite curved ellipse.

\clearpage
\subsection{Curvature and validation scope}
For our geometry, the local metric and total-angle measures are different:
\begin{equation}
 \frac{a}{R}=0.0025,\qquad \frac{L}{R}=2.5\ {\rm rad},\qquad
 \frac{\gamma\sigma_s}{\sigma_r}\simeq577,
 \quad \sigma_s=\frac{L}{\sqrt{12}},\quad\sigma_r\simeq\frac a2.
 \label{eq:literature-geometry}
\end{equation}
Small $a/R$ is consistent with a local weak-curvature expansion. The
$143.24^\circ$ span does not by itself invalidate it, while global Cartesian
slicing would misrepresent the arc. The RMS aspect ratio uses the radial
width consistently; the still larger ratio based on the thin vertical width
does not remove the radial or end-field question.

Qiang's 2025 solver derives all kicks from one potential and treats curvature
perturbatively for smooth long bunches. Its Gaussian aspect-ratio illustration
above 1000 is not a strict 1\% validity bound~\cite{qiang2025}.
The improved 2026 preprint adds free-space reference-radius matching in
Eqs. 48--54~\cite{qiang2026}; that material must not be attributed to the
published 2025 version.

The 2026 regularized Gaussian curvature potential adds a linear term to
remove an integral divergence. Its displayed correction has zero on-axis
radial derivative as written~\cite{qiang2026}; this derivative statement is
our inference. It therefore does not directly predict our finite-arc centre
field of $+0.055092$ V/m. Comparison requires matching the global physical
source, boundary conditions and reference-force convention, which is distinct
from numerical mesh padding. No intent to remove a reference force is inferred.

Zhou et al.'s HB2025 study shows why long-arc frame reconstruction must be
checked independently of the field closure~\cite{zhou2025}.
The accepted 2026 paper describes a dynamic orbit-based frame, but only its
abstract was accessible; no radius-specific error bound can be extracted
from it~\cite{zhou2026}. Our frame controls address this numerical concern;
they do not supply the missing nonlocal end or curvature field.

Validation observables also differ. ImpactX's Gaussian single-kick example
tests a reduced-model kick against the same model's analytic prediction;
its acceptance tolerances are not a universal 1\% physical guarantee~\cite{impactx}.
SACR instead supplies an independent Coulomb reference for our prescribed
physical density. The sign error at the reference-path midpoint and the
5 mm end-field deficit limit the accuracy claim for the archived closure in
this stress case. They neither invalidate smooth-bunch applications nor
establish tracking errors.

Our 43-target metrics use raw laboratory electric fields. A charge-weighted
transverse average for the curved source would instead be
\begin{equation}
 \overline E_s(s)=\frac{\int_A\rho J E_s\,\dd u\,\dd y}
 {\int_A\rho J\,\dd u\,\dd y},\qquad J=1+u/R.
 \label{eq:literature-average}
\end{equation}
This average has not been computed; the 43 diagnostic points are not a
charge-weighted quadrature.

These findings motivated the nonuniform-density SACR extension and nonlocal
OPEN operator in \cref{sec:free-open}. Its triangular reference uses a weighted
boundary plus density-gradient volume term, rather than the earlier uniform
boundary-only identity. The new test compares the same finite physical source
with its full Coulomb field. Varying-curvature paths, retarded fields and
self-consistent curved tracking remain separate validation questions.
